% !TEX root = ../../main.tex
\section{Mô tả bài toán}
\label{sec:mo_ta_bai_toan}

Trong hoạt động nghiên cứu khoa học tại các trường đại học và viện nghiên cứu hàng đầu nói chung, cũng như tại VinUniversity nói riêng, quá trình sản xuất các công bố học thuật quốc tế chất lượng cao (các hội nghị chuyên ngành danh giá như NeurIPS, ICML, CVPR và các tạp chí uy tín thuộc danh mục Q1/Q2) đòi hỏi học giả phải xử lý khối lượng tri thức khổng lồ với độ chính xác và tính liêm chính học thuật (\textit{Academic Integrity}) tuyệt đối.

Trước đây, khi tiến hành một công trình nghiên cứu mới, học giả, giảng viên và sinh viên thường phải thực hiện thủ công toàn bộ các công đoạn nghiên cứu thông qua nhiều công cụ đơn lẻ, rời rạc:
\begin{enumerate}[label=\textbf{(\arabic*)}]
    \item \textbf{Tìm kiếm và thu thập ấn phẩm học thuật:} Học giả phải chủ động truy cập độc lập vào từng cơ sở dữ liệu hoặc cổng lưu trữ như Google Scholar, arXiv, DBLP, ACM Digital Library hay IEEE Xplore. Việc tải các tệp PDF tài liệu về máy tính cá nhân diễn ra phân tán, dễ dẫn đến tình trạng lưu trữ trùng lặp các tệp dung lượng lớn, gây lãng phí tài nguyên và khó khăn trong việc quản trị phiên bản hay đồng bộ thư viện nghiên cứu.
    \item \textbf{Đọc hiểu và trích xuất dữ liệu thực nghiệm:} Khi mở tệp PDF trên các trình đọc tài liệu truyền thống hoặc sử dụng các công cụ trò chuyện với tài liệu phổ biến, các công cụ này thường chỉ bóc tách văn bản thô theo dòng quét một cột, làm xáo trộn hoàn toàn cấu trúc hai cột (\textit{two-column format}) đặc thù của các bài báo khoa học. Đáng lo ngại hơn, các mô hình ngôn ngữ lớn phổ quát hiện nay thường xuyên mắc phải \textbf{hiện tượng ảo giác trích dẫn (Citation Hallucination)}: tự bịa đặt các số liệu đo lường không có thật, hoặc chỉ trích dẫn số trang chung chung. Học giả phải mất hàng giờ đồng hồ rà soát thủ công từng dòng chữ, bảng biểu trên ấn phẩm gốc để kiểm chứng tính xác thực của thông tin.
    \item \textbf{Tổng hợp mạng lưới phương pháp và bộ dữ liệu:} Để có cái nhìn toàn cảnh về một lĩnh vực (State-of-the-Art), nhà nghiên cứu phải tự lập các bảng tính Excel hoặc ghi chú Notion rời rạc nhằm theo dõi xem phương pháp nào đã được thử nghiệm trên bộ dữ liệu nào và đạt kết quả ra sao. Cách làm thủ công này rất dễ bỏ sót các mối liên hệ bắc cầu giữa các công trình nghiên cứu và không hỗ trợ truy vấn so sánh liên tài liệu.
    \item \textbf{Soạn thảo bản thảo khoa học bằng \LaTeX:} Học giả phải chuyển ngữ cảnh sang các trình soạn thảo như TeXstudio hay Overleaf. Khi biên dịch thất bại, học giả phải tự lần mò qua các tệp nhật ký lỗi (\texttt{.log}) dài hàng nghìn dòng để phát hiện và sửa các lỗi cú pháp vi mô (thiếu thẻ đóng môi trường toán học, ký tự đặc biệt chưa thoát mã, hoặc nhãn trích dẫn bị xung đột).
\end{enumerate}

Sự chia cắt giữa các công cụ khiến người nghiên cứu mất trung bình từ $30\%$ đến $40\%$ tổng thời gian cho các thao tác chuyển đổi ngữ cảnh và xử lý kỹ thuật thủ công, thay vì tập trung vào tư duy sáng tạo khoa học cốt lõi.

Do đó, việc nghiên cứu và xây dựng nền tảng trợ lý AI nghiên cứu khoa học toàn trình \textbf{\VS} là hết sức cấp thiết. Nền tảng được thiết kế nhằm giải quyết triệt để bài toán ảo giác trích dẫn thông qua kiến trúc Agentic RAG lai (Hybrid Search RRF), tự động định vị khung viền tọa độ Bounding Box minh chứng trực tiếp trên bản PDF gốc, kết nối mạng tri thức học thuật 4 thực thể (GraphRAG), và tích hợp tác nhân AI tự phục hồi lỗi biên dịch \LaTeX\ trong môi trường máy ảo cách ly an toàn.
